\newpage
\phantomsection
\label{prf:monotone-composition-rules}
\LRAProofFor{thm:monotone-composition-rules}

\begin{remark*}[Return]
\hyperref[thm:monotone-composition-rules]{Return to Theorem}
\end{remark*}

\begin{theorem*}[Composition Rules for Monotonicity]
Order-preserving and order-reversing maps compose according to the parity of
the number of reversals.
\end{theorem*}

\begin{proof}[Professional Standard Proof]
\LRAProofBodyStart
Let $a_1\leq a_2$. If both maps preserve order, then
$f(a_1)\leq f(a_2)$ and $g(f(a_1))\leq g(f(a_2))$. If both reverse order,
then $f(a_2)\leq f(a_1)$ and applying $g$ reverses again, giving
$g(f(a_1))\leq g(f(a_2))$. If exactly one map reverses order, the final
inequality is reversed.
\end{proof}

\begin{proof}[Detailed Learning Proof]
\LRAProofBodyStart
Track the direction of the inequality through the two maps. A preserving map
keeps the direction. A reversing map flips it. Two flips cancel; one flip
remains.
\end{proof}
\begin{remark*}[Proof structure]
\end{remark*}



\begin{dependencies}
\begin{itemize}
  \item \hyperref[thm:monotone-composition-rules]{Proof target}
\end{itemize}
\end{dependencies}

\clearpage

